\section{Simulation-Free Training of Flows}
\label{sec:method}

%
%
%
We consider generative models that define a mapping between samples $x_1$ from a
noise distribution $p_1$ to samples $x_0$ from a data distribution $p_0$ in
terms of an ordinary differential equation (ODE),
\begin{equation}
  \label{eq:ode}
  dy_t = v_\Theta(y_t, t)\,dt\;,
\end{equation}
where the velocity $v$ is parameterized by the weights $\Theta$ of a neural network.
Prior work by~\citet{Chen2018NeuralOD} suggested to directly solve \Cref{eq:ode} via differentiable ODE solvers. However, this process is computationally expensive, especially for large network architectures that parameterize $v_\Theta(y_t, t)$. 
A more efficient alternative is to directly regress a vector field $u_t$ that generates a probability path between $p_0$ and $p_1$.
%
To construct such a $u_t$, we define a forward process, corresponding to a probability path $p_t$ between $p_0$ and $p_1=\mathcal{N}(0, 1)$, as
%
%
\begin{equation}
  \label{eq:forwardprocess}
  z_t = a_t x_0 + b_t \epsilon\quad\text{where}\;\epsilon \sim \mathcal{N}(0,I)\;. 
\end{equation}
For $a_0 = 1, b_0 = 0, a_1 = 0$ and $b_1 = 1$, the marginals,
\begin{align}
  \label{eq:marginals}
  p_t(z_t) &=
  \mathbb{E}_{\epsilon \sim \mathcal{N}(0,I)}
  p_t(z_t \vert \epsilon)\;,
\end{align}
are consistent with the data and noise distribution.

To express the relationship between $z_t, x_0$ and $\epsilon$, we introduce
%
$\psi_t$ and $u_t$ as 
\begin{align}
      \psi_t(\cdot | \epsilon) &: x_0 \mapsto a_t x_0 + b_t \epsilon \\
        u_t(z| \epsilon) &\coloneqq \psi'_t(\psi_t^{-1}(z| \epsilon)  \vert \epsilon) \label{eq:utmapsdiffinverse}
\end{align}

%
%
%
%
%
%
  %
%
%
%
%
Since $z_t$ can be written as solution to the ODE  $z_t' = u_t(z_t |\epsilon)$, with initial value $z_0=x_0$, $u_t(\cdot | \epsilon)$ generates $p_t(\cdot | \epsilon)$.
 %
Remarkably, one can construct a marginal vector field $u_t$ which generates the marginal probability paths $p_t$  \cite{lipman2023flow} (see \ref{subsec:flowproofs}), using the conditional vector fields $u_t(\cdot | \epsilon)$:
\begin{align}
    u_t(z) = \mathbb{E}_{\epsilon \sim
  \mathcal{N}(0,I)} u_t(z \vert \epsilon) \frac{p_t(z \vert
  \epsilon)}{p_t(z)}
  \label{eq:marginal_u}
\end{align}
While regressing $u_t$ with the \emph{Flow Matching} objective 
\begin{align}
   \mathcal{L}_{FM} =  \mathbb{E}_{t, p_t(z)} || v_{\Theta}(z, t) - u_t(z) ||_2^2.
\end{align}
directly is intractable due to the marginalization in \eqref{eq:marginal_u}, 
\emph{Conditional Flow Matching} (see \ref{subsec:flowproofs}),
\begin{align}
   \mathcal{L}_{CFM} =  \mathbb{E}_{t, p_t(z | \epsilon), p(\epsilon) }|| v_{\Theta}(z, t) - u_t(z | \epsilon)  ||_2^2\;, \label{eq:condflowmatch}
\end{align}
with the conditional vector fields $u_t(z \vert \epsilon)$ provides an equivalent yet tractable objective.

%
%
%

To convert the loss into an explicit form we insert
 $\psi_t'(x_0 \vert \epsilon) = a_t'x_0 + b_t' \epsilon$ and $\psi_t^{-1}(z\vert \epsilon) = \frac{z -b_t \epsilon}{a_t}$ into (\ref{eq:utmapsdiffinverse})
\begin{equation}
  z_t' = u_t(z_t \vert \epsilon) = \frac{a_t'}{a_t} z_t - \epsilon b_t (\frac{a_t'}{a_t} - \frac{b_t'}{b_t})\;.
  \label{eq:zode}
\end{equation}
%
%
%
%
%
%
%
Now, consider the \emph{signal-to-noise ratio} $\lambda_t := \log \frac{a_t^2}{b_t^2}$. 
With $\lambda_t' = 2 (\frac{a_t'}{a_t} - \frac{b_t'}{b_t})$, we can rewrite \Cref{eq:zode} as 
\begin{equation}
  u_t(z_t \vert \epsilon) = \frac{a_t'}{a_t} z_t - \frac{b_t}{2} \lambda_t' \epsilon 
  \label{eq:zodewithlambda}
\end{equation}

Next, we use \Cref{eq:zodewithlambda} to reparameterize \Cref{eq:condflowmatch} as a noise-prediction objective:
\begin{align}
    %
    \mathcal{L}_{CFM} &= \mathbb{E}_{t, p_t(z | \epsilon), p(\epsilon) } || v_{\Theta}(z, t) - \frac{a_t'}{a_t} z  + \frac{b_t}{2} \lambda_t' \epsilon ||_2^2 \\
    &= \mathbb{E}_{t, p_t(z | \epsilon), p(\epsilon) } \left(-\frac{b_t}{2}\lambda_t' \right)^2  || \epsilon_\Theta(z, t) - \epsilon ||_2^2
    %
    %
    %
\end{align}
%
where we defined $\epsilon_\Theta \coloneqq \frac{-2}{\lambda_t' b_t} (v_\Theta - \frac{a_t'}{a_t} z)$. 

Note that the optimum of the above objective does not change when
introducing a time-dependent weighting. Thus, one can derive various weighted
loss functions that provide a signal towards the desired solution but might
affect the optimization trajectory. For a unified analysis of different
approaches, including classic diffusion formulations, we can write the objective
in the following form (following \citet{Kingma2023UnderstandingDO}):
\begin{equation*}
  \mathcal{L}_w(x_0) = -\frac{1}{2} \mathbb{E}_{t\sim\mathcal{U}(t), \epsilon\sim \mathcal{N}(0, I)}
  \left[ w_t \lambda_t' \Vert \epsilon_\Theta(z_t, t) - \epsilon
  \Vert^2 \right]\;,
\end{equation*}
%
where $w_t = -\frac{1}{2} \lambda_t' b_t^2$ corresponds to $\mathcal{L}_{CFM}$.


%
%
%
%
%
%
%
%
%

%
%
%
%
%
%
%
%
%
%
%
%

%
%
%
%
%
%
%
%
%
%
%
%
%
%
%
%
%
%
%
%
%
%
%
%
%
%
%
%
%
%
%
%
%
%
%
%
%
%

%
%
%
%
%
%
%
%
%
%
%
%
%
%
%
  
%
%
%
%
%
%
%
%
%
%
%
%
%
%
%
%
%
%
%
%
%
%
%

%


%
%
\section{Flow Trajectories}
%
\label{subsec:variants}
%
%
%
%
In this work, we consider different variants of the above formalism that we
briefly describe in the following.

\paragraph{Rectified Flow}
Rectified Flows (RFs) \cite{liu2022flow,albergo2022building,lipman2023flow} define the forward process as straight paths
between the data distribution and a standard normal distribution, i.e.
\begin{equation}
z_t = (1-t) x_0 + t \epsilon\;,
\end{equation}
and uses $\mathcal{L}_{CFM}$ which then corresponds to $w_t^\text{RF} = \frac{t}{1-t}$. The network output directly parameterizes the velocity $v_\Theta$.
%
%
%
%
%
%
%
%
%


\paragraph{EDM}
EDM \cite{Karras2022ElucidatingTD} uses a forward process of the form
\begin{equation}
  z_t = x_0 + b_t \epsilon
\end{equation}
where \cite{Kingma2023UnderstandingDO}
%
$ b_t = \exp{F_{\mathcal{N}}^{-1}(t \vert P_m, P_s^2)} $
%
with $F_{\mathcal{N}}^{-1}$ being the quantile function of the normal
distribution with mean $P_m$ and variance $P_s^2$. Note that this choice results in
\begin{equation}
  \lambda_t \sim \mathcal{N}(-2P_m, (2P_s)^2)\quad\text{for}\;t\sim
  \mathcal{U}(0,1)
\end{equation}
The %
network is parameterized through an \textbf{F}-prediction
\cite{Kingma2023UnderstandingDO,Karras2022ElucidatingTD} and the loss can be
written as $\mathcal{L}_{w_t^\text{EDM}}$ with
\begin{equation}
  w_t^\text{EDM} = \mathcal{N}(\lambda_t \vert -2P_m, (2P_s)^2)(e^{-\lambda_t}+0.5^2)
\end{equation}

\paragraph{Cosine}
\cite{nichol2021improved} proposed a forward process of the form
\begin{equation}
  z_t = \cos\bigl(\frac{\pi}{2} t\bigr) x_0 + \sin\bigl(\frac{\pi}{2} t\bigr) \epsilon\;.
\end{equation}
In combination with an $\epsilon$-parameterization and loss, this corresponds
to a weighting $w_t = \operatorname{sech}(\lambda_t/2)$.
When combined with a \textbf{v}-prediction loss
\cite{Kingma2023UnderstandingDO}, the weighting is given by $w_t = e^{-\lambda_t/2}$.

\paragraph{(LDM-)Linear}
LDM \cite{Rombach_2022} uses a modification of the DDPM schedule \cite{ho2020denoising}. Both are
variance preserving schedules, i.e. $b_t=\sqrt{1-a_t^2}$, and define $a_t$ for
discrete timesteps $t = 0, \dots, T-1$ in terms of diffusion coefficients
$\beta_t$ as
%
$ a_t  = (\prod_{s=0}^t (1 - \beta_s))^{\frac{1}{2}}$.
%
For given boundary values $\beta_0$ and $\beta_{T-1}$, DDPM uses
%
$ \beta_t = \beta_0 + \frac{t}{T-1} (\beta_{T-1} - \beta_0) $
%
and LDM uses
%
$\beta_t = \left( \sqrt{\beta_0\vphantom{\beta_{T-1}}} + \frac{t}{T-1} (\sqrt{\beta_{T-1}} - \sqrt{\beta_0\vphantom{\beta_{T-1}}}) \right)^2$. 
%

\subsection{Tailored SNR Samplers for RF models}
\label{subsubsec:tailoredsnrforrf}
The RF loss trains the velocity $v_\Theta$ uniformly on all timesteps in
$[0, 1]$.
Intuitively, however, the resulting velocity prediction target $\epsilon - x_0$ is more
difficult for $t$ in the middle of $[0, 1]$, since for $t=0$, the optimal
prediction is the mean of $p_1$, and for $t=1$ the optimal prediction is the
mean of $p_0$. 
In general, changing the distribution over $t$ from the commonly used
uniform distribution $\mathcal{U}(t)$ to a distribution with density $\pi(t)$
is equivalent to a weighted loss $\mathcal{L}_{w_t^\pi}$ with
\begin{equation}
  w_t^\pi = \frac{t}{1-t}\pi(t)
\end{equation}
Thus, we aim to give more weight to intermediate timesteps by sampling them more frequently. Next, we describe the timestep densities $\pi(t)$ that we use to train our models.

%
%

\paragraph{Logit-Normal Sampling}
One option for a distribution that puts more weight on intermediate steps
is the logit-normal distribution~\cite{atchison1980logistic}. Its density,
\begin{equation}
\pi_{\text{ln}}(t; m, s) = \frac{1}{s\sqrt{2\pi}} \frac{1}{t(1-t)}\exp\Bigl(-\frac{(\text{logit}(t)-m)^2}{2s^2}\Bigr),
\label{eq:logitnormal}
\end{equation}
where $\text{logit}(t) = \log \frac{t}{1-t}$, has a location parameter, $m$,
and a scale parameter, $s$. The location parameter enables us to bias
the training timesteps towards either data $p_0$ (negative $m$) or noise $p_1$
(positive $m$). As shown in \Cref{fig:tsamplers}, the scale parameters controls how
wide the distribution is.

In practice, we sample the random variable $u$ from a normal distribution $u \sim \mathcal{N}(u; m, s)$ and map it through the standard logistic function.

\paragraph{Mode Sampling with Heavy Tails}
%
%
The logit-normal density always vanishes at the endpoints $0$ and $1$. To study
whether this has adverse effects on the performance, we also use a timestep
sampling distribution with strictly positive density on $[0, 1]$. For a scale
parameter $s$, we define
\begin{equation}
f_{\text{mode}}(u; s) = 1 - u - s\cdot\Bigl(\cos^2\bigl(\frac{\pi}{2} u\bigr) - 1 + u\Bigr).
\label{eq:modesampler}
\end{equation}
For $ -1 \leq s \leq \frac{2}{\pi -2}$, this function is
monotonic, and we can use it to sample from the implied density
$\pi_{\text{mode}}(t; s) = \left| \frac{d}{dt} f_{\text{mode}}^{-1}(t) \right|$.
%
%
As seen in \Cref{fig:tsamplers}, the scale parameter controls the degree to which either
the midpoint (positive $s$) or the endpoints (negative $s$) are favored during sampling.
This formulation also includes a \textbf{uniform weighting}
$\pi_{\text{mode}}(t; s=0) = \mathcal{U}(t)$ for $s=0$, which has been used
widely in previous works on Rectified Flows \cite{liu2022flow,ma2024sit}.

\paragraph{CosMap}
Finally, we also consider the \emph{cosine} schedule \cite{nichol2021improved}
from \Cref{subsec:variants} in the RF setting. 
In particular, we are looking for a mapping $f: u \mapsto f(u) = t, \; u \in [0, 1]$, %
such that the log-snr matches that of the cosine schedule: $2 \log \frac{\cos(\frac{\pi}{2}u)}{\sin(\frac{\pi}{2}u)} = 2 \log \frac{1-f(u)}{f(u)}$. 
Solving for $f$, we obtain for $u \sim \mathcal{U}(u)$
\begin{equation}
t = f(u) = 1 - \frac{1}{\tan(\frac{\pi}{2} u) + 1}, 
\label{eq:cosmap}
\end{equation}
from which we obtain the density 
\begin{equation}
\pi_{\text{CosMap}}(t) = \left| \frac{d}{dt} f^{-1}(t) \right| = \frac{2}{\pi - 2\pi t + 2\pi t^2}.
\label{eq:cosmapdensity}
\end{equation}
%
%



%
%




%
%
%
%
%

\section{Text-to-Image Architecture}
\label{subsec:methodarch}
\modelarchitecture
%
\vspace{-0.5em}
For text-conditional sampling of images, our model has to take both modalities,
text and images, into account. We use pretrained models to derive suitable
representations and then describe the architecture of our diffusion backbone.
An overview of this is presented in \cref{fig:modelfig:total}.

Our general setup follows LDM \citep{Rombach_2022} for training text-to-image models in the latent space of a pretrained autoencoder.
Similar to the encoding of images to latent representations, we also follow
previous approaches \cite{saharia2022photorealistic, balaji2022ediffi} and encode the text conditioning $c$ using
pretrained, frozen text models. Details can be found in \Cref{suppsec:detailsonreps}. 

%
%
%
%
%
%
%
%
%
%
%
%
%

%
%
%
%
%
%
%
%
%
%
%
%
%
%
%
%
%
%

\textbf{Multimodal Diffusion Backbone}
%
%
%
%
%
%
%
%
Our architecture builds upon the DiT \cite{Peebles_2023}
architecture. DiT only considers class conditional image generation and uses a
modulation mechanism to condition the network on both the timestep of the
diffusion process and the class label. Similarly, we use
embeddings of the timestep $t$ and $c_\text{vec}$ as inputs to the
modulation mechanism. However, as the pooled text representation retains only
coarse-grained information about the text input \cite{podell2023sdxl}, the network also
requires information from the sequence representation $c_\text{ctxt}$.

We construct a sequence consisting of embeddings of the text and image inputs.
Specifically, we add positional encodings and flatten $2\times 2$ patches of
the latent pixel representation $x \in \RR^{h \times w \times c}$ to a patch
encoding sequence of length $\frac{1}{2} \cdot h \cdot \frac{1}{2} \cdot w$. 
After embedding this
patch encoding and the text encoding $c_\text{ctxt}$ to a common
dimensionality, we concatenate the two sequences. We then follow DiT and apply
a sequence of modulated attention and MLPs. %
%

Since text and image embeddings are conceptually quite different, we use two
separate sets of weights for the two modalities. As shown in \cref{fig:modelfig:mm}, this
is equivalent to having two independent transformers for each modality, but
joining the sequences of the two modalities for the attention operation, such
that both representations can work in their own space yet take the other one
into account.

For our scaling experiments, we parameterize the size of the model in terms of
the model's depth $d$, \ie the number of attention blocks, by setting the
hidden size to $64\cdot d$ (expanded to $4\cdot64\cdot d$ channels in the MLP
blocks), and the number of attention heads equal to $d$.
%


